% GENERATED FILE. Edit canonical lessons, then run npm run generate:books.
% canonical-source-hash: fnv1a64:f02fe27327c8beb1
% canonical-lessons: GE-C385-ausnahme, GE-C385-bestrafen, GE-C385-gesetzlich, GE-C385-rechtlich, GE-C385-regelung, GE-R385-review-378-385

\chapter{Regeln und Ausnahmen}
\label{ch:ge-regeln-und-ausnahmen}
\hlchaptermodality{385}

\begin{chapteropening}
\textbf{By the end of this chapter:} \emph{I can talk about rules and exceptions, what the law lays down, and punishment.}

Complete the last lesson of chapter 385: I can talk about rules and exceptions, what the law lays down, and punishment.
\end{chapteropening}

\section[die Ausnahme]{die Ausnahme — the exception}

\label{lesson:GE-C385-ausnahme}



\begin{quote}
Before the new one: say the German for to extend, then the German for to expire.
\end{quote}

\subsection*{You'll want to know: die Ausnahme}
\textbf{die Ausnahme} — \textquotedblleft{}the exception\textquotedblright{}.

\begin{quote}
Normalerweise arbeite ich nicht am Sonntag, aber heute ist eine Ausnahme.
\end{quote}

\begin{quote}
\emph{Normally I do not work on Sunday, but today is an exception.}
\end{quote}

Plural: \emph{die Ausnahmen}. \emph{Mit Ausnahme von} — with the exception of.

\begin{cousinweb}
\emph{Ausnahme} is the noun of \emph{ausnehmen} \textquotedblleft{}to take out, to except\textquotedblright{}, built on \emph{nehmen} \textquotedblleft{}to take\textquotedblright{}, Proto-Germanic \emph{nemaną} (EtymDB): a case taken out of the rule.
\end{cousinweb}

\begin{grammarlens}[title={What you've built}]
One more word for telling, explaining and describing at B1.
\end{grammarlens}

\subsection*{Your turn}
\begin{itemize}
\raggedright
  \item \emph{Say it:} \emph{die Ausnahme}
  \item \emph{Say it:} \emph{die Ausnahme}, once more
  \item \emph{Say it:} a sentence of your own with \emph{die Ausnahme}
  \item \emph{Recall:} say the German for to extend, then the German for to expire, then say \emph{die Ausnahme} again
\end{itemize}

\begin{tcolorbox}[breakable,colback=teal!4,colframe=teal!35!black,title={Before you move on}]
What does die Ausnahme mean? (\textquotedblleft{}The exception\textquotedblright{}.) Say it once more.
\end{tcolorbox}
\section[bestrafen]{bestrafen — to punish}

\label{lesson:GE-C385-bestrafen}



\begin{quote}
Before the new one: say the German for to expire, then the German for the exception.
\end{quote}

\subsection*{You'll want to know: bestrafen}
\textbf{bestrafen} — \textquotedblleft{}to punish\textquotedblright{}.

\begin{quote}
Der Schüler wurde für die Lüge bestraft.
\end{quote}

\begin{quote}
\emph{The pupil was punished for the lie.}
\end{quote}

Inseparable: \emph{er bestraft, er hat bestraft}. \emph{Die Strafe} — the punishment, also a fine.

\begin{cousinweb}
\emph{bestrafen} is built with \emph{be-} on \emph{strafen} \textquotedblleft{}to punish\textquotedblright{}, inherited from Middle High German \emph{strāfen} (EtymDB); EtymDB lists Dutch \emph{straffen} as a cousin.
\end{cousinweb}

\begin{grammarlens}[title={What you've built}]
One more word for telling, explaining and describing at B1.
\end{grammarlens}

\subsection*{Your turn}
\begin{itemize}
\raggedright
  \item \emph{Say it:} \emph{bestrafen}
  \item \emph{Say it:} \emph{bestrafen}, once more
  \item \emph{Say it:} a sentence of your own with \emph{bestrafen}
  \item \emph{Recall:} say the German for to expire, then the German for the exception, then say \emph{bestrafen} again
\end{itemize}

\begin{tcolorbox}[breakable,colback=teal!4,colframe=teal!35!black,title={Before you move on}]
What does bestrafen mean? (\textquotedblleft{}To punish\textquotedblright{}.) Say it once more.
\end{tcolorbox}
\section[gesetzlich]{gesetzlich — legal, laid down by law}

\label{lesson:GE-C385-gesetzlich}



\begin{quote}
Before the new one: say the German for the exception, then the German for to punish.
\end{quote}

\subsection*{You'll want to know: gesetzlich}
\textbf{gesetzlich} — \textquotedblleft{}legal, laid down by law\textquotedblright{}.

\begin{quote}
Der Mindestlohn ist gesetzlich geregelt.
\end{quote}

\begin{quote}
\emph{The minimum wage is laid down by law.}
\end{quote}

\emph{Gesetzlich verboten} — prohibited by law. \emph{Der gesetzliche Feiertag} — the public holiday.

\begin{cousinweb}
\emph{gesetzlich} is built with \emph{-lich} on \emph{Gesetz} \textquotedblleft{}law\textquotedblright{}, which EtymDB derives from \emph{setzen} \textquotedblleft{}to set\textquotedblright{}, Proto-Germanic \emph{satjaną}: what has been set down.
\end{cousinweb}

\begin{grammarlens}[title={What you've built}]
One more word for telling, explaining and describing at B1.
\end{grammarlens}

\subsection*{Your turn}
\begin{itemize}
\raggedright
  \item \emph{Say it:} \emph{gesetzlich}
  \item \emph{Say it:} \emph{gesetzlich}, once more
  \item \emph{Say it:} a sentence of your own with \emph{gesetzlich}
  \item \emph{Recall:} say the German for the exception, then the German for to punish, then say \emph{gesetzlich} again
\end{itemize}

\begin{tcolorbox}[breakable,colback=teal!4,colframe=teal!35!black,title={Before you move on}]
What does gesetzlich mean? (\textquotedblleft{}Legal, laid down by law\textquotedblright{}.) Say it once more.
\end{tcolorbox}
\section[rechtlich]{rechtlich — legally, in law}

\label{lesson:GE-C385-rechtlich}



\begin{quote}
Before the new one: say the German for to punish, then the German for legal.
\end{quote}

\subsection*{You'll want to know: rechtlich}
\textbf{rechtlich} — \textquotedblleft{}legally, in law\textquotedblright{}.

\begin{quote}
Rechtlich ist das kein Problem.
\end{quote}

\begin{quote}
\emph{Legally, that is no problem.}
\end{quote}

\emph{Rechtlich} is about the law in general; \emph{gesetzlich} about what a statute says.

\begin{cousinweb}
\emph{rechtlich} is built with \emph{-lich} on \emph{Recht} \textquotedblleft{}law, right\textquotedblright{}, inherited from Old High German \emph{reht} and Proto-Germanic \emph{rehtaz} \textquotedblleft{}right, straight\textquotedblright{} (EtymDB), a cousin of English \emph{right}.
\end{cousinweb}

\begin{grammarlens}[title={What you've built}]
One more word for telling, explaining and describing at B1.
\end{grammarlens}

\subsection*{Your turn}
\begin{itemize}
\raggedright
  \item \emph{Say it:} \emph{rechtlich}
  \item \emph{Say it:} \emph{rechtlich}, once more
  \item \emph{Say it:} a sentence of your own with \emph{rechtlich}
  \item \emph{Recall:} say the German for to punish, then the German for legal, then say \emph{rechtlich} again
\end{itemize}

\begin{tcolorbox}[breakable,colback=teal!4,colframe=teal!35!black,title={Before you move on}]
What does rechtlich mean? (\textquotedblleft{}Legally, in law\textquotedblright{}.) Say it once more.
\end{tcolorbox}
\section[die Regelung]{die Regelung — the rule (that applies), the arrangement}

\label{lesson:GE-C385-regelung}



\begin{quote}
Before the new one: say the German for legal, then the German for legally.
\end{quote}

\subsection*{You'll want to know: die Regelung}
\textbf{die Regelung} — \textquotedblleft{}the rule (that applies), the arrangement\textquotedblright{}.

\begin{quote}
Für Kinder unter sechs gibt es eine besondere Regelung.
\end{quote}

\begin{quote}
\emph{There is a special rule for children under six.}
\end{quote}

Plural: \emph{die Regelungen}. From \emph{regeln} — to regulate, to sort out.

\begin{cousinweb}
\emph{Regelung} is the noun of \emph{regeln}, built on \emph{Regel} \textquotedblleft{}rule\textquotedblright{}, from Old High German \emph{regula} and Latin \emph{rēgula} \textquotedblleft{}a straight piece of wood, a ruler, a rule\textquotedblright{} (EtymDB; Lewis \& Short).
\end{cousinweb}

\begin{grammarlens}[title={What you've built}]
One more word for telling, explaining and describing at B1.
\end{grammarlens}

\subsection*{Your turn}
\begin{itemize}
\raggedright
  \item \emph{Say it:} \emph{die Regelung}
  \item \emph{Say it:} \emph{die Regelung}, once more
  \item \emph{Say it:} a sentence of your own with \emph{die Regelung}
  \item \emph{Recall:} say the German for legal, then the German for legally, then say \emph{die Regelung} again
\end{itemize}

\begin{tcolorbox}[breakable,colback=teal!4,colframe=teal!35!black,title={Before you move on}]
What does die Regelung mean? (\textquotedblleft{}The rule (that applies), the arrangement\textquotedblright{}.) Say it once more.
\end{tcolorbox}
\section[(dialogue)]{Words from chapters 378–385 — a review}

\label{lesson:GE-R385-review-378-385}



\begin{quote}
Say the German for the term paper and the class talk.
\end{quote}

\subsection*{Your turn}
Say these aloud:

\begin{itemize}
\raggedright
  \item the term paper, the class talk, the essay, the summary, to sum up
  \item to rate, the evaluation, to grade, to make up, to correct
  \item the admission, to let in, the scholarship, the registration, to withdraw
  \item the regulation, to lay down, to observe, to keep to, obliged
  \item the ban, permitted, punishable, to break, the violation
  \item the authority, the application, to apply for, the form, the signature
  \item the permit, to approve, valid, to extend, to expire
  \item the exception, to punish, legal, legally, the rule
\end{itemize}

\begin{tcolorbox}[breakable,colback=teal!4,colframe=teal!35!black,title={Before you move on}]
The term paper? (\textbf{die Hausarbeit}.) The class talk? (\textbf{das Referat}.) The essay? (\textbf{der Aufsatz}.) The summary? (\textbf{die Zusammenfassung}.) To sum up? (\textbf{zusammenfassen}.) To rate? (\textbf{bewerten}.) The evaluation? (\textbf{die Bewertung}.) To grade? (\textbf{benoten}.) To make up? (\textbf{nachholen}.) To correct? (\textbf{korrigieren}.) The admission? (\textbf{die Zulassung}.) To let in? (\textbf{zulassen}.) The scholarship? (\textbf{das Stipendium}.) The registration? (\textbf{die Anmeldung}.) To withdraw? (\textbf{abmelden}.) The regulation? (\textbf{die Vorschrift}.) To lay down? (\textbf{vorschreiben}.) To observe? (\textbf{beachten}.) To keep to? (\textbf{einhalten}.) Obliged? (\textbf{verpflichtet}.) The ban? (\textbf{das Verbot}.) Permitted? (\textbf{zulässig}.) Punishable? (\textbf{strafbar}.) To break? (\textbf{verstoßen}.) The violation? (\textbf{der Verstoß}.) The authority? (\textbf{die Behörde}.) The application? (\textbf{der Antrag}.) To apply for? (\textbf{beantragen}.) The form? (\textbf{das Formular}.) The signature? (\textbf{die Unterschrift}.) The permit? (\textbf{die Genehmigung}.) To approve? (\textbf{genehmigen}.) Valid? (\textbf{gültig}.) To extend? (\textbf{verlängern}.) To expire? (\textbf{ablaufen}.) The exception? (\textbf{die Ausnahme}.) To punish? (\textbf{bestrafen}.) Legal? (\textbf{gesetzlich}.) Legally? (\textbf{rechtlich}.) The rule? (\textbf{die Regelung}.)
\end{tcolorbox}
